\section{Introduction}
\label{sec:intro}

\suppressfloats[t]
\begin{figure}[b!]
\centering
\includegraphics[width=\linewidth]{Figures/goalscope/goalscope-teaser-with-eval.png}
\caption{\textbf{The GoalScope readout setup.} (a) At the highlighted pivot (\texttt{`` b"}), we cache the selected attention heads' outputs that are added to the residual stream. (b) We aggregate these outputs into a single vector $v_g$. (c) In a separate LM run, the user may ask different queries, such as (c.1) asking the LM to verbalize the next step, (c.2) implement the next step in a different programming language. We steer the LM response by adding $v_g$ at each generation position. The query can be adapted to elicit different aspects of the model's encoding of its context and upcoming behavior. The right panel shows that we can get decent readouts as early as the indent (\texttt{``    "}) token, 3 tokens before the concrete implementation begins.}
\label{fig:teaser}
\label{fig:goalscope-setup}
\end{figure}


\arnab{I am using terms like ``sub-goal'', ``intention'', ... interchangeably in the intro. We need to come to a consensus on what term to use.}
\lisa{I like subgoal/goal}

Consider a programmer asked to implement \emph{merge sort}, a task that decomposes into several sub-tasks: \emph{split the array}, \emph{sort each half}, and \emph{merge the sorted halves}. Before writing the code for the final step, the programmer already has a coherent goal: \textit{``merge the two halves.''} We ask whether large language models (LLMs) write code the same way: when an LLM emits a long span of tokens to satisfy a multi-step specification, does it hold an explicit representation of the \emph{next} instrumental sub-goal before it begins its implementation?

\lisa{could we use planning instead of instrumental? otherwise maybe we should use instrumental in the abstract}

We study the \emph{instrumental} level of this hierarchy: the \emph{broader goal} (implement merge sort) decomposes into \emph{instrumental sub-goals} (merge the two sorted arrays), each of which admits many \emph{implementations} (such as using a \emph{for} loop or a \emph{while} loop). Our question concerns a representation of \emph{which sub-goal the model is about to pursue}, and whether that representation can be separated from its concrete realizations in code.

Although LLMs are trained to predict one token at a time, a growing body of work shows that they \emph{look ahead}: they encode information about tokens they have not yet emitted \citep{pal2023futurelens} and plan toward constraints several tokens in the future \citep{lindsey2025biology}. We ask what the plan for the \emph{next step} looks like as a concrete internal object, and whether that object is one we can read and move between contexts.\arnab{maybe move to Related Works if the intro gets long.}
\lisa{this could also be moved up to support the context?}

We investigate this question through activation patching on code tasks with multiple sub-goals. Our main experiments use \texttt{Qwen2.5-Coder-32B-Instruct}, with replications across model families and sizes. Our analysis yields three key findings.

\textbf{The next sub-goal is encoded before implementation, and is verbalizable (\Cref{sec:what}).} We localize the encoding of the LLM's next sub-goal to a compact set of attention heads, which we call the \goalheads. We introduce \goalscope, a lightweight probe built on these heads that verbalizes the upcoming sub-goal directly from the LM's internal activations. Before the model writes the implementation, we can extract what the heads write to the residual stream and use it to generate a faithful description of the operation it is about to perform.

\textbf{The encoding is abstract and portable (\Cref{sec:what:readout}).} Transplanting the encoding steers a fresh code context to pursue the source sub-goal. A goal extracted from Python can trigger the same operation in C++, and vice versa. This portability provides evidence that the heads carry a representation of the operation that generalizes across its surface realizations. We also explore whether the same heads can verbalize sub-goals in code whose component steps are not explicitly listed, and inferred relations in non-code few-shot tasks (\Cref{sec:generalization}).

\textbf{Goal heads mediate the model's implementation of the sub-goal (\Cref{sec:heads,sec:mechanism}).} A small set of heads suffices for effective readout, and a broader set is causally important for carrying out the operation. When we restrict other upper-layer heads' access to the upcoming task and code span, preserving goal-head access largely preserves implementation performance. Conversely, restricting the broader goal-head set to the attention sink sharply reduces performance, while restricting random heads has little effect. The heads attend to relevant task descriptions, and changing their query or key/value states can redirect the sub-goal pursued. Using their value and output projections, we also introduce \textbf{OVReader} to trace where the upcoming sub-goal becomes readable in intermediate residual states.

Our results provide concrete, testable evidence that code LLMs represent and pursue intermediate goals: a compact, portable, causally manipulable encoding of \emph{intent} that precedes action.
\arnab{Teaser figure: the readout + one steering example, framed as ``intent before action''.}
